% Third test document for todofloats.sty: one-column text with side=alternate,
% where the notes of each page take turns between the right and the left
% margin, and leftmarginsep keeps the left notes clear of a strip next to the
% text (line numbers, in a review copy).  One note in the left margin is long
% enough to be cut, so its continuation must keep the same gap.  check.sh
% compiles it and checks the layout and the page size.
\documentclass[11pt]{article}
\usepackage[a4paper,margin=2.5cm,heightrounded=true]{geometry}
\usepackage{hyperref}
\usepackage{lipsum}
\usepackage{xcolor}
\newcommand{\todofontsize}{\fontsize{7}{8}\selectfont}
\usepackage[textsize=todofontsize,grow=40mm,side=alternate,leftmarginsep=15mm,authors={timv=magenta!10},agents={claude}]{mytodonotes} % [disable]

\begin{document}
\section{Alternating Margins}
\lipsum[1]
A first note\timv{A1 goes to the right margin.} here,
a second\claude{A2 goes to the left margin, 15mm from the text.} close to it,
and a third\timv{A3 goes to the right margin again, below A1.} on the same line.
\lipsum[2]
A long note\timv{A4 is in the left margin and long enough to be cut. \lipsum[3-7]} here.
\lipsum[8-16]
On the second page\claude{A5 is the first note of its page, so it goes right.} the turns start again,
and\timv{A6 is the second note of its page: left.} so on.
\lipsum[17-18]
\begin{figure}[h]
  \centering\fbox{a figure}
  \caption{A caption longer than one line, so that LaTeX typesets it twice: once
    to measure it, then for real.  The note in it\timv{A7 sits in a two-line
    caption; its measuring copy never reaches the page.} must not make every
    run ask for another one.}
\end{figure}
\lipsum[19]
\end{document}
